This Online Appendix details the specifications behind the spillover tests summarized in Section~\ref{sec:spillovers}.

The baseline test estimates the effect of \emph{Plan Cuadrante} on reported crime, arrests, police officers, and patrol vehicles in neighboring untreated municipalities,
\begin{equation}
Outcome_{mt}= \delta_{1}\,\mbox{\textit{Plan Cuadrante near municip.}}_{mt} + Municipality_{m} + Year_{t} + u_{mt}
\label{eq:main3}
\end{equation}
\noindent where the sample is restricted to municipalities that never adopted the program, each is assigned a treatment date equal to the earliest adoption among its treated neighbors, and $\delta_{1}$ is estimated dynamically with the staggered difference-in-differences estimator of \citet{Callaway}. Because these administrative outcomes are observed for every municipality, we estimate the specification on the full national sample of never treated municipalities (Online Appendix Figure~\ref{fig:spillover_admin}).

For the intensity-based test in the spirit of \citet{crepon2013}, we replace the binary neighbor-treatment indicator with a continuous measure of exposure and restrict the sample to never-treated municipalities,
\begin{equation}
Outcome_{mt}= \delta_{1}\,\mbox{\textit{Share neighbors with Plan Cuadrante}}_{mt} + Municipality_{m} + Year_{t} + u_{mt}
\label{eq:crepon}
\end{equation}
\noindent Because the treatment is continuous, we estimate it with the heterogeneity-robust estimator of \citet{Chaisemartin}, measuring exposure both as the share of treated neighbors (Online Appendix Figure~\ref{fig:figure_continuous_spillovers}) and as their number (Online Appendix Figure~\ref{fig:figure_count_spillovers}).

For the randomization-inference test, we reassign adoption dates and never-treated status across municipalities $1{,}000$ times and re-estimate the direct and spillover effects under the sharp null, yielding the placebo distributions in Online Appendix Figure~\ref{fig:ri_placebo}. For this exercise we rely on the estimator of \citet{Chaisemartin} rather than that of \citet{Callaway}: re-estimating the latter over a thousand permutations would require prohibitive computational power and running time.

Finally, we re-estimate the baseline administrative specification on the municipalities neighboring the 46 ENUSC municipalities (Online Appendix Figure~\ref{fig:figure_46ENUSC_spillovers_admin}); and, for the outcomes measured only in ENUSC, we estimate the spillover of early-adopting ENUSC municipalities on their late-adopting ENUSC neighbors before the latter's own adoption (Online Appendix Figure~\ref{fig:figure_results_spillovers_enusc}). The six focal late-adopters, each shown with its earliest-adopting ENUSC neighbor and that neighbor's adoption year, are Molina (Curic\'o, 2005), Tocopilla (Iquique, 2005), San Carlos (Chill\'an, 2007), Padre Hurtado (Pe\~naflor, 2008), Calera (Quillota, 2008), and Limache (Quillota, 2008). Across all of these specifications, the estimated spillover effects are small and statistically indistinguishable from zero.
